Let \(\mathcal T_v\) be the finite signed-isotone response-table set and put \(\mathcal T=\prod_{v\in V}\mathcal T_v\).  A complete table vector \(t\in\mathcal T\), together with the acyclic directed order of \(G\), determines every potential-response literal in an unnested event by recursive substitution.  Thus, for \(E=A_Q\cap B_Q\), there is a well-defined indicator \(e_E(t)\in\{0,1\}\).

For a district \(D\), set
\[
 C_D(t_D)=\{(R_v:v\in D)=t_D\},
 \qquad
 C(t)=\bigcap_{D\in\mathcal D(G)}C_D(t_D).
\]
The exact response-table atoms are pairwise disjoint and partition the formal complete response-table space.  Consequently,
\[
 \mathbf 1_E
 =\sum_{t\in\mathcal T:e_E(t)=1}\mathbf 1_{C(t)}
 =\sum_{t\in\mathcal T:e_E(t)=1}
   \prod_{D\in\mathcal D(G)}\mathbf 1_{C_D(t_D)}.
\]
This is a finite signed sum, with every nonzero coefficient equal to \(+1\), and each summand has exactly one response-type cell per district.

The expansion is unique in the fixed exact-atom basis of Definition \(\mathrm{def:order\text{-}normal\text{-}form}\).  Indeed, if \(\mathbf 1_E=\sum_{t\in\mathcal T}c_t\mathbf 1_{C(t)}\) as a formal atom expansion, evaluation at the atom indexed by \(t_0\) gives \(c_{t_0}=e_E(t_0)\).  The stipulated lexicographic ordering changes only the order of display.

For every signed parent coordinate, Assumption \(\mathrm{ass:signed\text{-}local\text{-}monotonicity}\) makes the normalized response table isotone.  A cell demanding output one at \(z\) and output zero at a comparable \(z'\ge z\) is therefore empty.  These are exactly the local restrictions deleted by the order normalizer, so every deleted cell has probability zero under every \(\mathfrak m\in\mathfrak M(G,M)\).

Every exogenous index affecting a vertex in a district \(D\) has \(C(u)=D\).  Hence distinct district response vectors depend on disjoint exogenous-coordinate blocks.  The product exogenous law in \(\mathrm{ass:semi\text{-}markovian}\) gives
\[
 P_{\mathfrak m}\{C(t)\}
 =\prod_{D\in\mathcal D(G)}P_{\mathfrak m}\{C_D(t_D)\}.
\]
Taking expectations in the atom identity yields
\[
 P_{\mathfrak m}(A_Q\cap B_Q)
 =\sum_{t:e_E(t)=1}
   \prod_{D\in\mathcal D(G)}P_{\mathfrak m}\{C_D(t_D)\},
\]
which is the value of \(\mathsf N_{G,M}(Q;A_Q\cap B_Q)\).  Repeating the same argument with \(E=B_Q\) proves the second assertion.